\chapter{Fundamentação Teórica}
\label{cap:fundamentacao}

Este capítulo apresenta a base conceitual que sustenta o desenvolvimento e a arquitetura do sistema de apoio ao Barema do COLCIC. São abordados os princípios de arquitetura de sistemas para a Web, os conceitos de processamento de documentos binários e mineração de texto, além dos fundamentos de \textit{interfaces} reativas e persistência de dados.

%\section{Conceitos Fundamentais}


\section{Arquitetura em Camadas e o Padrão MVC}

O desenvolvimento de aplicações web modernas exige a adoção de modelos arquiteturais que favoreçam a modularidade e o isolamento de responsabilidades. Tradicionalmente, o padrão arquitetural em camadas organiza o sistema em níveis distintos — apresentação, negócio e persistência —, de forma que alterações em uma camada não comprometam as demais. Dentro desse modelo,  o padrão \textit{Model-View-Controller} (MVC) é amplamente empregado para separar a lógica de apresentação da lógica de negócio e de acesso a dados \cite{pressman2016}.

Fowler descreve o MVC como um dos padrões arquiteturais mais relevantes para aplicações web, destacando que a separação entre o modelo de dados, a camada de visualização e o controlador de fluxo é o que permite a 
evolução independente de cada componente sem ruptura do sistema como um todo \cite{fowler2002}. No contexto deste trabalho, essa separação é crítica: as regras de negócio do COLCIC precisam estar isoladas da 
interface e da camada de geração de documentos, de modo que alterações no regulamento possam ser incorporadas sem reescrever o sistema inteiro.


\section{Mapeamento Objeto-Relacional e Persistência de Dados}

A camada de persistência de dados em sistemas orientados a objetos frequentemente enfrenta o problema da incompatibilidade de impedância, caracterizado pela diferença entre o modelo conceitual de objetos na memória e o modelo relacional das tabelas em um banco de dados SQL. Para mitigar esse problema, utiliza-se o Mapeamento Objeto-Relacional (\textit{Object-Relational Mapping} — ORM).

O ORM atua como uma camada de abstração que traduz de forma transparente as operações sobre objetos computacionais em comandos de manipulação de dados em bancos relacionais \cite{fowler2002}. Essa abstração permite que a aplicação opere sobre estruturas nativas da linguagem — como dicionários e listas de objetos em Python — e delegue ao ORM a responsabilidade de sincronizar essas estruturas com o banco de dados, 
minimizando o acoplamento direto com o motor de persistência escolhido.


\section{Processamento de Documentos Digitais e Extração de Texto}

A manipulação de arquivos no formato PDF (\textit{Portable Document Format}) introduz desafios relacionados à natureza não estruturada das informações. Diferente de linguagens de marcação estruturadas (como XML ou JSON), o PDF foca na fidelidade de renderização visual e na disposição física de caracteres em coordenadas bidimensionais, omitindo frequentemente tags de estrutura semântica ou separadores lógicos de parágrafos.

Para transformar esses fluxos de bytes brutos em dados inteligíveis pela lógica de negócio, utilizam-se bibliotecas especializadas na decodificação da estrutura interna do arquivo. A partir do conteúdo textual extraído, aplicam-se técnicas de processamento baseadas em Expressões Regulares (\textit{Regular Expressions} — Regex). As expressões regulares operam como autômatos finitos determinísticos que realizam o casamento de padrões textuais \cite{hopcroft2006}, permitindo varrer o conteúdo do documento em busca de tokens numéricos adjacentes a qualificadores de carga horária — como "horas", "h" e "carga horária" — para identificar automaticamente os metadados operacionais dos certificados


\section{Reatividade de Interface e Tempo de Resposta}

A experiência do usuário em plataformas utilitárias é diretamente influenciada pela latência observada nas atualizações da interface. Nielsen estabelece três limiares fundamentais de tempo de resposta que permanecem 
relevantes desde suas formulações originais: respostas em até 0,1 segundo sustentam a percepção de instantaneidade; respostas em até 1,0 segundo mantêm o fluxo de pensamento do usuário sem interrupções perceptíveis; e respostas que ultrapassam 10 segundos comprometem a atenção do usuário e exigem \textit{feedback} visual explícito de progresso \cite{nielsen1993}.

Para alcançar esse patamar de reatividade em operações de manipulação de arquivos, adota-se o modelo de computação assíncrona na camada de apresentação. Utilizando \textit{APIs} nativas dos navegadores modernos, como a \textit{Fetch API}, a interface web envia dados e arquivos binários em segundo plano para as rotas do servidor. O servidor processa as informações e devolve respostas leves em formato JSON, permitindo que o estado visual do navegador seja modificado dinamicamente via JavaScript sem a necessidade de recarregamento completo da página e sem interromper o fluxo de pensamento do usuário.

%\section{Tecnologias e Ferramentas}

%\subsection{Microframeworks Web: Flask}

%Em sistemas que operam com processamento transiente de dados, a sobrecarga de frameworks monolíticos pode comprometer a eficiência do servidor. Nesse cenário, destacam-se os microframeworks, que fornecem apenas o  núcleo essencial para o roteamento HTTP e o gerenciamento de requisições, delegando funcionalidades adicionais a extensões acopladas sob demanda. O Flask é um microframework Python que segue esse princípio, oferecendo ao desenvolvedor controle direto sobre a estrutura da aplicação sem impor convenções rígidas \cite{grinberg2018}. Essa flexibilidade permite estruturar um orquestrador baseado em serviços na memória da aplicação, viabilizando o gerenciamento de estados voláteis sem dependências rígidas de componentes globais.

%\subsection{Biblioteca de Processamento de PDF: PyMuPDF}

%O PyMuPDF é uma biblioteca Python que fornece \textit{bindings} para a engine MuPDF, permitindo o acesso programático à estrutura interna de documentos PDF \cite{pymupdf2026}. A biblioteca possibilita a extração do número de páginas, do conteúdo textual e de metadados do documento diretamente a partir do fluxo de bytes do arquivo em memória, sem necessidade de escrita temporária em disco. Essa característica é especialmente relevante para o contexto deste trabalho, onde os certificados precisam ser inspecionados de forma imediata no momento do upload, antes mesmo de serem persistidos permanentemente no servidor.